% ch07_u02 —— 原书 p.246–p.255（第 7 章 7.3.4–7.4.10）
% 底本：A. Beilinson, V. Drinfeld, Quantization of Hitchin's Integrable System
% and Hecke Eigensheaves，附录「代数栈上的 D-模理论与 Hecke 模式」

\SourcePage{246}

以 $D(\mathcal{Y})^{\ge0}\subset D(\mathcal{Y})$ 记由下述 Ω-复形 $F$ 构成的满子范畴：当 $a<0$ 时 $H^{a}_{D}F=0$。类似地定义 $D(\mathcal{Y})^{\le0}$。

\bdp{7.3.4} \emph{命题.} 这给出 $D(\mathcal{Y})$ 上的一个 t-结构，其核心为 $\mathcal{M}(\mathcal{Y})$，上同调函子为 $H_{\mathcal{D}}$。

此命题由下文的引理 7.5.3 立即推出。当 $\mathcal{Y}$ 是轨道栈时，另有一个证明，见 7.6.11。

\bdp{7.3.5} \emph{注.} 考虑函子 $\Omega:C(\mathcal{M}(\mathcal{Y}))\to C(\mathcal{Y},\Omega)$。对 $M\in C(\mathcal{M}(\mathcal{Y}))$ 有 $H^{\cdot}M=H_{\mathcal{D}}(\Omega M)$，故 Ω 给出 t-正合函子 $\Omega:D(\mathcal{M}(\mathcal{Y}))\to D(\mathcal{Y})$，它延拓两个核心之间的「恒等」等价。若 $\mathcal{Y}$ 是 Deligne--Mumford 栈\footnote{即存在一个簇，它平展地覆盖 $\mathcal{Y}$。此时函子 $D:C(\mathcal{Y},\Omega)\to C(\mathcal{M}(\mathcal{Y}))$ 显然有意义（它给出逆等价 $D(\mathcal{M}(\mathcal{Y}))\to D(\mathcal{Y})$，见 7.2.7）。}，则此函子是范畴等价；但一般情形下不然。

\bdp{7.3.6} 设 $f:\mathcal{Y}\to\mathcal{Z}$ 是满足（310）的光滑栈的态射。它给出 DG 赋环拓扑的态射 $(\mathcal{Y}_{sm},\Omega_{\mathcal{Y}})\to(\mathcal{Z}_{sm},\Omega_{\mathcal{Z}})$，从而给出一对伴随 DG 函子
\begin{equation}\tag{311}
f_{\Omega}^{\cdot}:C(\mathcal{Z},\Omega)\to C(\mathcal{Y},\Omega),\qquad
f_{\cdot}:C(\mathcal{Y},\Omega)\to C(\mathcal{Z},\Omega),
\end{equation}
以及同伦范畴之间相应的伴随三角函子（因 $\mathcal{Y}$ 拟紧，$f_{\cdot}$ 保持拟凝聚）。

若 $f$ 光滑，则 $f_{\Omega}^{\cdot}$ 保持 $\mathcal{D}$-拟同构，故它定义 t-正合函子 $f^{\cdot}:D(\mathcal{Z})\to D(\mathcal{Y})$。显然它与 $f$ 的复合相容。
% [Tutor Note: 原书 p.246 此句印作 "f_Ω preservres D-quasi-isomorphisms"，"preservres" 系 "preserves" 之误，底本如此，保留原样。]

设 $f$ 任意。把推动函子 $f_{*}:D^{+}(\mathcal{Y})\to D^{+}(\mathcal{Z})$ 定义为右导出函子 $Rf_{\cdot}$。我们将在 7.3.10 证明 $f_{*}$ 定义正确。为此需要足够多的「脆弱」对象。

\SourcePage{247}

\bdp{7.3.7} \emph{定义.} 称 $\mathcal{Y}$ 上的 $\mathcal{O}$-模 $F$ 为松的，若对 $\mathcal{Y}$ 上任意平坦 $\mathcal{O}$-模 $P$ 都有 $H^{a}(\mathcal{Y},P\otimes F)=0$（$a>0$）。$\mathcal{O}$-复形或 Ω-复形 $F$ 为松的，若每个 $F^{i}$ 都松。

\bdp{7.3.8} \emph{引理.} (i) 对 $\mathcal{Y}$ 上任意 Ω-复形 $F'$，存在 $\mathcal{D}$-拟同构 $F'\to F$，且 $F$ 松。若 $F'$ 下有界，则可取 $F$ 下有界。(ii) 设 $f$（见 7.3.6）光滑且仿射。则 $f_{\Omega}^{\cdot}$ 与 $f_{\cdot}$ 把松 Ω-复形送到松 Ω-复形。(iii) 若 $F_{1},F_{2}$ 分别是栈 $\mathcal{Y}_{1},\mathcal{Y}_{2}$ 上的松 Ω-复形，则 $F_{1}\boxtimes F_{2}$ 是 $\mathcal{Y}_{1}\times\mathcal{Y}_{2}$ 上的松 Ω-复形。

\emph{证明.} (i) 因 $\mathcal{Y}$ 拟紧，存在 $\mathcal{Y}$ 的超覆叠 $U_{\cdot}$，其中各 $U_{a}$ 是仿射概形。因 $\mathcal{Y}$ 的对角态射仿射，投影 $\pi_{a}:U_{a}\to\mathcal{Y}$ 也仿射。取 $F$ 为 $F'$ 关于此超覆叠的 Čech 复形，即
\[
F^{i}=\bigoplus_{a\ge0}\pi_{a\cdot}\bigl(F'^{i-a}_{U_{a}}\bigr).
\]
(ii) 显然。(iii) 可设 $F_{i}$ 是松 $\mathcal{O}_{\mathcal{Y}_{i}}$-模。设 $P$ 是 $\mathcal{Y}_{1}\times\mathcal{Y}_{2}$ 上的平坦 $\mathcal{O}$-模。因 $F_{1}$ 松，对 $a>0$ 有 $R^{a}p_{2\cdot}(P\otimes p_{1}^{*}F_{1})=0$，且 $p_{2\cdot}(P\otimes p_{1}^{*}F_{1})$ 是 $\mathcal{Y}_{2}$ 上的平坦 $\mathcal{O}$-模（这里 $p_{i}:\mathcal{Y}_{1}\times\mathcal{Y}_{2}\to\mathcal{Y}_{i}$ 是投影）。于是
\[
H^{a}\bigl(\mathcal{Y}_{1}\times\mathcal{Y}_{2},P\otimes(F_{1}\boxtimes F_{2})\bigr)
=H^{a}\bigl(\mathcal{Y}_{2},\bigl(p_{2\cdot}(P\otimes p_{1}^{*}F_{1})\bigr)\otimes F_{2}\bigr),
\]
因 $F_{2}$ 松，此式当 $a>0$ 时为零。\hfill$\Box$

回到 7.3.6 末尾的情形。

\bdp{7.3.9} \emph{引理.} 若 $F$ 是 $\mathcal{Y}$ 上下有界的松 Ω-复形，则 $f_{\cdot}F=Rf_{\cdot}F$。

\emph{证明.} 只需验证：若 $F$ 还是 $\mathcal{D}$-无环的（即满足 $H^{\cdot}_{D}F=0$），则 $f_{\cdot}F$ 也 $\mathcal{D}$-无环（用 7.3.8(i)）。

a. 可设 $\mathcal{Z}$ 是光滑仿射概形 $Z$。事实上，待验证的陈述关于 $\mathcal{Z}$ 是局部的。把 $\mathcal{Z}$ 换成 $\mathcal{Z}_{sm}$ 中的仿射对象 $Z$，把 $\mathcal{Y}$ 换成 $\mathcal{Y}\times_{\mathcal{Z}}Z$，
\SourcePage{248}%
并把 $F$ 换成它到 $\mathcal{Y}\times_{\mathcal{Z}}Z$ 的拉回。新数据满足引理的全部条件。

b. 可设 $\mathcal{Y}$ 是光滑仿射概形 $Y$。事实上，取 (i) 中的 $U_{\cdot}$，记 $A$ 为项是 $A^{i}=\bigoplus_{a\ge0}(f\pi_{a})_{\cdot}(F'^{i-a}_{U_{a}})$ 的 Čech 复形。$A$ 是 $\mathcal{Z}$ 上的 Ω-复形。因 $F$ 松，显然的态射 $f_{\cdot}F\to A$ 是 $\mathcal{D}$-拟同构（用（310））。注意 $A$ 带显然的滤链，其相继商为 $(f\pi_{a})_{\cdot}(F_{U_{a}})[-a]$。若这些复形都 $\mathcal{D}$-无环，则 $A$ 也 $\mathcal{D}$-无环（此处用 $F$ 下有界），故 $f_{\cdot}F$ 是 $\mathcal{D}$-无环的。

c. 设 $i:Y\to Y\times Z$ 是 $f$ 的图像嵌入。则 $G:=i_{\cdot}F$ 是 $\mathcal{D}$-无环的。因 $f_{\cdot}F=p_{\cdot}G$（$p$ 是投影 $Y\times Z\to Z$），只需证 $p_{\cdot}G$ 是 $\mathcal{D}$-无环的。设 $T$ 是 $\mathcal{D}G$ 沿 $p$ 的纤维的相对 de Rham 复形。处于直积情形，故 $p_{\cdot}T$ 是 $Z$ 上的 $\mathcal{D}$-复形。有显然的 $\mathcal{D}$-复形态射 $Dp_{\cdot}G\to p_{\cdot}T$，它是拟同构。因 $p_{\cdot}T$ 无环（$T$ 带滤链，其相继商为 $DG\otimes\Lambda\Theta_{Y}$，而 $\mathcal{D}G$ 无环），证毕。\hfill$\Box$

\emph{注.} 设 $f$ 是仿射态射。则对任意 $F\in C(\mathcal{Y},\Omega)$ 有 $f_{\cdot}F=Rf_{\cdot}F$。事实上，该陈述关于 $\mathcal{Z}$ 是局部的，故可设 $\mathcal{Z}$ 是仿射概形。此时 $\mathcal{Y}$ 是仿射概形，于是 $\mathcal{Y}$ 上任意复形都松；再用 7.3.9 即可。

\bdp{7.3.10} \emph{推论.} (i) 函子 $f_{*}:=Rf_{\cdot}:D^{+}(\mathcal{Y})\to D^{+}(\mathcal{Z})$ 定义正确。(ii) $f_{*}$ 与 $f$ 的复合相容，即典范态射 $(f_{1}f_{2})_{*}\to f_{1*}f_{2*}$ 是同构。

\emph{证明.} (i) 用 7.3.8(i) 与 7.3.9。(ii) $f_{\cdot}$ 把松 Ω-复形送到松 Ω-复形。\hfill$\Box$

\bdp{7.3.11} \emph{注.} (i) 上述引理在 $\mathcal{O}$-复形的框架中也成立。

\SourcePage{249}
(ii) 设 $\mathcal{Y}$ 的 $\mathcal{O}$-模范畴上的函子 $f_{\cdot}$ 有有限的上同调维数（例如 $f$ 可表时如此）。则对任意边界条件的 Ω-复形的导出范畴，$f_{*}:=Rf_{\cdot}$ 都定义正确。事实上，7.3.9（连同其证明）对无界的松 Ω-复形仍成立。

(iii) 若我们的栈是光滑簇，则上述函子 $f_{*}$ 就是 $\mathcal{D}$-模论中标准的推动函子（见 7.2.10）。此时引理 7.3.9（连同其证明）仍成立，只需假定：对 $Y$ 的任意使 $U\to Y$ 为仿射态射的 Zariski 开集 $U$，上同调 $H^{a}(U,F^{i})$（$a>0$）为零。

\bdp{7.3.12} 现在设 $\mathcal{Y}$ 是任一光滑栈，其对对角态射 $\mathcal{Y}\to\mathcal{Y}\times\mathcal{Y}$ 仿射（即放弃拟紧性假设）。此时 $\mathcal{Y}$ 上 Ω-复形的范畴可能太小，不足以定义正确的 $\mathcal{D}$-模导出范畴。我们如下扩充上述形式体系。

为简化记号，设 $\mathcal{Y}$ 有可数个拟紧开集构成的覆盖。换言之，$\mathcal{Y}$ 是递增的开拟紧子栈序列 $\mathcal{Y}_{1}\subset\mathcal{Y}_{2}\subset\ldots$ 之并。
% [Tutor Note: 原书 p.249 此处印作 "open quasi-comact substacks"，"quasi-comact" 系 "quasi-compact" 之误，底本如此，保留原样。]
$\mathcal{Y}_{\cdot}$ 上的 Ω-复形是集合 $F=(F_{i},a_{i})$，其中 $F_{i}$ 是 $\mathcal{Y}_{i}$ 上的 Ω-复形，而 $a_{i}:F_{i+1}|_{\mathcal{Y}_{i}}\to F_{i}$ 是 $\mathcal{D}$-拟同构。这样的 Ω-复形构成 DG 范畴 $C(\mathcal{Y}_{\cdot},\Omega)$，于是有相应的同伦范畴 $K(\mathcal{Y}_{\cdot},\Omega)$。它带有上同调函子 $H_{\mathcal{D}}$，取值于 $\mathcal{Y}$ 上 $\mathcal{D}$-模的阿贝尔范畴 $\mathcal{M}(\mathcal{Y})$：
\[
H_{\mathcal{D}}(F)|_{\mathcal{Y}_{i}}=H_{\mathcal{D}}(F_{i}).
\]

把 $D(\mathcal{Y}_{\cdot},\Omega)$ 定义为 $K(\mathcal{Y}_{\cdot},\Omega)$ 关于 $\mathcal{D}$-拟同构的局部化。不同 $\mathcal{Y}_{\cdot}$ 给出的三角范畴 $D(\mathcal{Y}_{\cdot},\Omega)$ 被典范地等同。事实上，设 $\mathcal{Y}'_{j}$ 是 $\mathcal{Y}$ 的另一个如上所述的开子栈序列。选取递增函数 $j=j(i)$，使 $\mathcal{Y}_{i}\subset\mathcal{Y}'_{j(i)}$。令 $\mathcal{Y}'_{\cdot}$ 上的 Ω-复形 $F'$ 对应 $\mathcal{Y}_{\cdot}$ 上的 Ω-复形 $F$，其中 $F^{i}=F'^{j(i)}|_{\mathcal{Y}_{i}}$。此函子与 $H_{\mathcal{D}}$ 交换。\SourcePage{250}%
导出范畴之间的相应函子在显然意义下不依赖辅助选取 $j(i)$，且是范畴等价。

我们看到，范畴 $D(\mathcal{Y}_{\cdot},\Omega)$ 只依赖于 $\mathcal{Y}$，故把它记作 $D(\mathcal{Y},\Omega)$ 或简记为 $D(\mathcal{Y})$。我们的三角范畴带有上同调函子 $H_{\mathcal{D}}:D(\mathcal{Y})\to\mathcal{M}(\mathcal{Y})$，并有典范的全忠实嵌入 $\Omega:\mathcal{M}(\mathcal{Y})\to D(\mathcal{Y})$（见 7.3.2）。命题 7.3.4 仍然成立；其证明由 7.5.4 推出。

设 $f:\mathcal{Y}\to\mathcal{Z}$ 是满足我们假设的光滑栈的态射。若 $f$ 光滑，则按显然方式定义 t-正合拉回函子 $f^{\cdot}:D(\mathcal{Z})\to D(\mathcal{Y})$。若 $f$ 是任意拟紧态射，则有典范推动函子 $f_{*}:D(\mathcal{Y})^{+}\to D(\mathcal{Z})^{+}$。我们在关于松 Ω-复形的简短插叙之后定义它。

按定义，$F\in C(\mathcal{Y}_{\cdot},\Omega)$ 是松的，若所有 $F_{i}\in C(\mathcal{Y}_{i},\Omega)$ 都松。引理 7.3.8(i),(iii) 在我们的情形仍成立。这意味着，可以只用松复形来定义 $\mathcal{D}$-导出范畴。为证 7.3.8(i)，取覆盖 $\pi_{i}:V_{i}\to\mathcal{Y}_{i}$，其中 $V_{i}$ 是仿射概形。以 $U_{i}$ 记诸 $V_{j}$（$1\le j\le i$）的不交并，以 $U_{i\cdot}$ 记 $\mathcal{Y}_{i}$ 的相应超覆叠，$U_{ia}$ 是 $U_{i}$ 在 $\mathcal{Y}_{i}$ 上的 $a$ 重纤维积。现取任意 $F'\in C(\mathcal{Y}_{\cdot},\Omega)$。设 $F_{i}$ 是 $F'_{i}$ 关于超覆叠 $U_{i\cdot}$ 的 Čech 复形（见 7.3.8(i) 的证明）。则诸 $F_{i}$ 按显然方式构成 $\mathcal{Y}_{\cdot}$ 上的 Ω-复形 $F$。此 $F$ 是松的，且显然的态射 $F'\to F$ 是 $\mathcal{D}$-拟同构，证毕。

现定义 $f_{*}$。设 $\mathcal{Z}_{i}$ 是 $\mathcal{Z}$ 的如上所述的开拟紧子栈序列。则 $\mathcal{Y}_{i}:=f^{-1}\mathcal{Z}_{i}$ 是 $\mathcal{Y}$ 的相应序列。设 $F$ 是 $\mathcal{Y}_{\cdot}$ 上下有界的松 Ω-复形。则 $(f_{\cdot}F)_{i}:=f_{\cdot}(F_{i})$ 构成 $\mathcal{Z}_{\cdot}$ 上的 Ω-复形 $f_{\cdot}F$（用 7.3.9）。函子 $f_{\cdot}$ 保持 $\mathcal{D}$-拟同构（由 7.3.9）。我们的 $f_{*}$ 就是 $\mathcal{D}$-导出范畴之间的相应函子。推论 7.3.10(ii) 连同其证明仍成立。

\SourcePage{251}

再设所有函子 $f_{i\cdot}:\mathcal{M}(\mathcal{Y}_{i},\mathcal{O})\to\mathcal{M}(\mathcal{Z}_{i},\mathcal{O})$ 都有有限的上同调维数（例如 $f$ 可表时如此）。则函子 $f_{*}$ 在整个 $D(\mathcal{Y})$ 上定义正确。事实上，设 $F$ 是 $\mathcal{Y}_{\cdot}$ 上任意松 Ω-复形。则 $(f_{\cdot}F)_{i}:=f_{\cdot}(F_{i})$ 构成 $\mathcal{Z}_{\cdot}$ 上的 Ω-复形 $f_{\cdot}F$（用 7.3.11(ii)）。函子 $f_{\cdot}$ 保持 $\mathcal{D}$-拟同构，故把 $f_{*}:D(\mathcal{Y})\to D(\mathcal{Z})$ 定义为 $\mathcal{D}$-导出范畴之间的相应函子。

\bdp{7.3.13} \emph{注.} 设 $A$ 是交换代数。以 $\mathcal{M}(\mathcal{Y},A)$ 记 $\mathcal{Y}$ 上带 $A$-作用的 $\mathcal{D}$-模构成的阿贝尔范畴。如 7.3.12 那样，用带 $A$-作用的 Ω-复形定义以 $\mathcal{M}(\mathcal{Y},A)$ 为核心的 t-范畴 $D(\mathcal{Y},A)$。标准的函子毫无困难地迁移到 $A$-线性的情形。更一般地，设 $A_{\mathcal{Y}}$ 是 $\mathcal{Y}$ 上的交换 $\mathcal{D}$-代数（即张量范畴 $\mathcal{M}^{\ell}(\mathcal{Y})$ 中的交换代数）。有 $A_{\mathcal{Y}}$-模的阿贝尔范畴 $\mathcal{M}(\mathcal{Y},A_{\mathcal{Y}})$ 及其导出版本 $D(\mathcal{Y},A_{\mathcal{Y}})$，后者如 7.3.12 那样用带 $A_{\mathcal{Y}}$-作用的 Ω-复形定义。

\bdsec{7.4. 导出范畴的下降}

我们解释导出范畴的一般同伦逆向极限构造。为能给出 $\mathcal{D}$-模导出范畴的「局部」定义，需要这一构造。

\bdp{7.4.1} 以 $(\Delta)$ 记非空有限全序集 $\Delta_{n}=[0,n]$ 与递增单射构成的范畴。设 $\mathcal{M}_{\cdot}$ 是余纤维化于 $(\Delta)$ 上的一族阿贝尔范畴，使得对任意态射 $\alpha:\Delta_{n}\hookrightarrow\Delta_{m}$，相应的函子 $\alpha_{\cdot}:\mathcal{M}_{n}\to\mathcal{M}_{m}$ 正合。

以 $\mathcal{M}_{\mathrm{tot}}$ 记 $\mathcal{M}_{\cdot}$ 的余笛卡儿截面构成的范畴：$\mathcal{M}_{\mathrm{tot}}$ 的对象是集合 $M=\{M_{n},\alpha^{*}\}$，其中 $M_{n}\in\mathcal{M}_{n}$，$\alpha^{*}=\alpha^{*}_{M}:\alpha_{\cdot}M_{n}\xrightarrow{\sim}M_{m}$ 是同构，且满足 $(\alpha\beta)^{*}=\alpha^{*}\alpha_{\cdot}(\beta^{*})$（这里 $\beta:\Delta_{l}\hookrightarrow\Delta_{n}$）。这是阿贝尔范畴。注意 $\mathcal{M}_{\mathrm{tot}}$ 与对偶相容：有 $(\mathcal{M}_{\mathrm{tot}})^{\circ}=(\mathcal{M}^{\circ})_{\mathrm{tot}}$。

\SourcePage{252}

我们的目标是定义以 $\mathcal{M}_{\mathrm{tot}}$ 为核心的 t-范畴 $D_{\mathrm{tot}}(\mathcal{M}_{\cdot})$，使其具有下述关键性质：
\begin{equation}\tag{312}
\begin{gathered}
\text{对任意 }M,N\in\mathcal{M}_{\mathrm{tot}}\text{，存在收敛到 }\operatorname{Ext}^{p+q}_{D_{\mathrm{tot}}(\mathcal{M}_{\cdot})}(N,M)\text{ 的典范谱序列 }E^{p,q}_{r}\text{，且}\\
E^{p,q}_{1}=\operatorname{Ext}^{q}_{\mathcal{M}_{p}}(N_{p},M_{p}).
\end{gathered}
\end{equation}
$D_{\mathrm{tot}}(\mathcal{M}_{\cdot})$ 的构造与对偶相容。

\bdp{7.4.2} 考虑范畴 $\sec_{+}=\sec_{+}(\mathcal{M}_{\cdot})$，其对象是集合 $M=(M_{n},\alpha^{*})$，其中 $M_{n}\in\mathcal{M}_{n}$，$\alpha^{*}=\alpha^{*}_{M}:\alpha_{\cdot}M_{n}\to M_{m}$ 是使 $(\alpha\beta)^{*}=\alpha^{*}\alpha_{\cdot}(\beta^{*})$、$\mathrm{id}^{*}_{\Delta_{n}}=\mathrm{id}_{M_{n}}$ 的态射。这是阿贝尔范畴，它包含 $\mathcal{M}_{\mathrm{tot}}$ 作为对扩张封闭的满子范畴。用对偶定义 $\sec_{-}=\sec_{-}(\mathcal{M}_{\cdot})$：$\sec_{-}(\mathcal{M}_{\cdot}):=(\sec_{+}(\mathcal{M}^{\circ}_{\cdot}))^{\circ}$；于是 $\sec_{-}$ 的对象是集合 $N=(N_{n},\alpha_{*})$，$N_{n}\in\mathcal{M}_{n}$，$\alpha_{*}=\alpha^{N}_{*}:N_{m}\to\alpha_{\cdot}N_{n}$。

考虑 $\sec_{\pm}$ 中复形构成的 DG 范畴 $C\sec_{\pm}$ 及相应的同伦范畴 $K\sec_{\pm}$。有伴随 DG 函子
\begin{equation}\tag{313}
c_{+}:C\sec_{-}\longrightarrow C\sec_{+},\qquad c_{-}:C\sec_{+}\longrightarrow C\sec_{-},
\end{equation}
它们定义如下。取 $M\in C\sec_{+}$。则对任意 $m\ge0$，有单纯形 $\Delta_{m}$ 上取值于 $C\mathcal{M}_{m}$ 的「上同调型」系数系统 $\widetilde{M}_{m}$。确切地说，$\widetilde{M}_{m}$ 把面 $\alpha:\Delta_{n}\hookrightarrow\Delta_{m}$ 对应到复形 $\alpha_{\cdot}M_{n}$；若 $\alpha':\Delta_{l}\hookrightarrow\Delta_{m}$ 是 $\alpha$ 的面，即 $\alpha'=\alpha\beta$，则相应的连接态射 $\alpha'_{\cdot}M_{l}\to\alpha_{\cdot}M_{n}$ 是 $\alpha_{\cdot}(\beta^{*})$。此时 $(c_{-}M)_{m}$ 是全上链复形 $C^{\cdot}(\Delta_{m},\widetilde{M}_{m})$（即 $c_{-}(M)^{\cdot}_{m}=\bigoplus_{\alpha:\Delta_{n}\to\Delta_{m}}\alpha_{\cdot}M_{n}^{\cdot-n}$），而 $\alpha^{c_{-}(M)}_{*}$ 是显然的投影。用对偶定义 $c_{+}$。

为看出 $c_{\pm}$ 伴随，对如上 $N,M$ 考虑阿贝尔群复形 $Hom(N,M)$，其项为
\[
Hom(N,M)^{i}=\prod_{a,n}\operatorname{Hom}\bigl(N^{a+n}_{n},M^{a+i}_{n}\bigr),
\]
而微分为 $df$，它把 $f=(f_{a,n})\in Hom(N,M)^{i}$ 送到
% [Tutor Note: 原书 p.252 此处印作 "the differential wich sends"，"wich" 系 "which" 之误，底本如此，保留原样。]
\[
(df)_{a,n}=df_{a,n}-(-1)^{i+n}f_{a+1,n}d+\sum_{j=0,\ldots,n}(-1)^{j}\alpha^{*}_{j}\alpha_{j\cdot}(f_{a+1,n-1})\alpha_{j*}.
\]

\SourcePage{253}

这里 $\alpha_{j}:\Delta_{n-1}\to\Delta_{n}$ 是第 $j$ 个面嵌入。于是伴随性由同态复形的显然等同
\begin{equation}\tag{314}
Hom(c_{+}N,M)\xrightarrow{\sim}Hom(N,M)\xleftarrow{\sim}Hom(N,c_{-}M)
\end{equation}
推出。

\bdp{7.4.3} \emph{注.} 固定某个 $m\ge0$。对 $i=0,\ldots,m$，设 $\nu_{i}:c_{-}(M)_{m}\to M_{m}$ 是投影 $c_{-}(M)_{m}\to\alpha_{i\cdot}M_{0}$ 与 $\alpha^{*}_{i}:\alpha_{i\cdot}M_{0}\to M_{m}$ 的复合；这里 $\alpha_{i}:\Delta_{0}\to\Delta_{m}$ 是第 $i$ 个顶点。此时诸态射 $\nu_{i}$ 彼此同伦（伴有典范同伦与「高阶同伦」）。

\bdp{7.4.4} \emph{引理.} 函子 $c_{\pm}$ 保持拟同构。伴随态射 $c_{+}c_{-}M\to M$、$N\to c_{-}c_{+}N$ 都是拟同构。\hfill$\Box$

由此可见，$c_{\pm}$ 给出导出范畴 $D\sec_{\pm}$ 之间互逆的等价。把如此等同的这两个范畴记作 $D\sec$。于是 $D\sec$ 带两个 t-结构，其核心分别为 $\sec_{\pm}$，上同调函子为 $H_{\pm}:D\sec\to\sec_{\pm}$。

\bdp{7.4.5} 设 $C_{\mathrm{tot}+}\subset C\sec_{+}$ 是由下述复形 $M$ 构成的满子范畴：对任意 $i$ 有 $H^{i}M\in\mathcal{M}_{\mathrm{tot}}\subset\sec_{+}$。换言之，$M\in C\sec_{+}$ 属于 $C_{\mathrm{tot}+}$，当且仅当所有态射 $\alpha^{*}_{M}$ 都是拟同构。类似地定义 $C_{\mathrm{tot}-}\subset C\sec_{-}$。以 $K_{\mathrm{tot}\pm}\subset K\sec_{\pm}$、$D_{\mathrm{tot}\pm}\subset D\sec_{\pm}$ 记相应的同伦范畴与导出范畴；它们是三角范畴。

导出范畴 $D(\mathcal{M}_{n})$ 构成 $(\Delta)$ 上的余纤维化范畴。以 $D^{\mathrm{fake}}_{\mathrm{tot}}$ 记其余笛卡儿截面构成的范畴（这\emph{不是}三角范畴！）。$\mathcal{M}_{\cdot}$ 的上同调函子给出函子 $H:D^{\mathrm{fake}}_{\mathrm{tot}}\to\mathcal{M}_{\mathrm{tot}}$。有显然的函子 $\epsilon_{+}:D_{\mathrm{tot}+}\to D^{\mathrm{fake}}_{\mathrm{tot}}$，它把 $M$ 对应到视为 $D^{\mathrm{fake}}_{\mathrm{tot}}$ 对象的资料 $(M_{n},\alpha_{*})$。类似地有函子 $\epsilon_{-}:D_{\mathrm{tot}}\to D^{\mathrm{fake}}_{\mathrm{tot}}$。

\bdp{7.4.6} \emph{引理.} 对任意 $M\in D_{\mathrm{tot}+}$ 有 $c_{-}M\in D_{\mathrm{tot}-}$，且存在唯一同构 $\epsilon_{-}(c_{-}M)\xrightarrow{\sim}\epsilon_{+}(M)$，其第 $0$ 个分量是 $\mathrm{id}_{M_{0}}$。把 $+$ 与 $-$ 互换，即得对偶的陈述。

\emph{证明.} 用 7.4.3。\hfill$\Box$

\SourcePage{254}

\bdp{7.4.7} 我们看到，函子 $c_{\pm}$ 等同三角范畴 $D_{\mathrm{tot}\pm}$。换言之，子范畴 $D_{\mathrm{tot}\pm}\subset D\sec$ 重合；这就是 7.4.1 中许诺的范畴 $D_{\mathrm{tot}}=D_{\mathrm{tot}}(\mathcal{M}_{\cdot})$。函子 $\epsilon_{\pm}$ 被典范地等同，故有函子 $\epsilon:D_{\mathrm{tot}}\to D^{\mathrm{fake}}_{\mathrm{tot}}$。注意 $H_{\pm}=H\epsilon$，故有典范上同调函子 $H:D_{\mathrm{tot}}\to\mathcal{M}_{\mathrm{tot}}$。这是 $D_{\mathrm{tot}}$ 上以 $\mathcal{M}_{\mathrm{tot}}$ 为核心的非退化 t-结构的上同调函子。注意嵌入 $D_{\mathrm{tot}}\to D\sec$ 关于 $D\sec$ 上两个 $\pm$ t-结构中的任一个都 t-正合；它把核心 $\mathcal{M}_{\mathrm{tot}}$ 与核心 $\sec_{+}$、$\sec_{-}$ 之交等同。

\bdp{7.4.8} 我们从 7.4.1 导出谱序列（312）。更一般地，考虑对象 $N\in D^{-}\sec_{-}\subset D\sec$、$M\in D^{+}\sec_{+}\subset D\sec$。用复形 $N\in K^{-}\sec_{-}$、$M\in K^{+}\sec_{+}$ 表示它们。考虑复形 $Hom(N,M)$（见 7.4.2）。它带显然的递减滤过 $F^{\cdot}$，且 $gr^{n}_{F}=Hom(N_{n},M_{n})[-n]$。注意 $Hom(N,M)$ 是下有界复形，且滤过 $F^{\cdot}$ 在每个项 $Hom(N,M)^{i}$ 上诱导有限滤过。我们把 $Hom(N,M)$ 视为此类复形的滤过导出范畴 DF 的对象。设 $RHom(N,\cdot)$ 是函子 $K^{+}\sec_{+}\to DF$、$M\mapsto Hom(N,M)$ 的右导出函子。此函子定义正确；且对任意 $n$，显然的态射 $gr^{n}_{F}RHom(N,M)\to RHom(N_{n},M_{n})[-n]$ 是拟同构。这由下述事实推出：对 $\mathcal{M}_{n}$ 中任意拟同构 $f:M_{n}\to I$，存在 $K^{+}\sec_{+}$ 中的拟同构 $g:M\to J$ 与态射 $h:I\to J_{n}$，使 $g_{n}=hf$。考虑滤过复形 $RHom(N,M)$ 的谱序列 $E^{p,q}_{r}$。它收敛到 $H^{\cdot}RHom(N,M)$，且
\[
E^{p,q}_{1}=H^{q}RHom_{\mathcal{M}_{p}}(N_{p},M_{p}).
\]

\bdp{7.4.9} \emph{注.} 设范畴 $\mathcal{M}_{n}$ 有许多内射对象。则范畴 $K^{+}_{\mathrm{tot}-}$ 有许多内射对象（即函子 $K^{+}_{\mathrm{tot}-}\to D^{+}_{\mathrm{tot}}$ 有右伴随函子）。事实上，若 $I\in K^{+}_{\mathrm{tot}+}$ 是使每个 $I^{a}_{n}$ 都是 $\mathcal{M}_{n}$ 的内射对象的复形，则 $c_{-}I$
% [Tutor Note: 原书 p.254 此处印作 "eacb"，系 "each" 之误，底本如此，保留原样。]
\SourcePage{255}%
是 $K^{+}_{\mathrm{tot}-}$ 的内射对象，而 $K^{+}_{\mathrm{tot}-}$ 中任意对象都拟同构于这样的 $I$。对偶地，若 $\mathcal{M}_{n}$ 有许多投射对象，则 $K^{-}_{\mathrm{tot}+}$ 有许多投射对象。

\bdp{7.4.10} 本小节后文不再使用；读者可以跳过。可以用稍许不同的方式定义 $D\sec$，从而定义 $D_{\mathrm{tot}}$，这在某些应用中较方便\footnote{此构造可追溯到 Toledo 与 Tong 的工作。}。如下定义范畴 $\mathrm{hot}_{+}=\mathrm{hot}_{+}(\mathcal{M}_{\cdot})$：其对象是族 $A=(A_{m})$，$A_{m}\in\mathcal{M}_{m}$。态射 $f:A\to B$ 是集合 $(f_{\alpha})$，其中对箭 $\alpha:\Delta_{n}\to\Delta_{m}$，相应的 $f_{\alpha}$ 是态射 $\alpha_{\cdot}A_{n}\to B_{m}$。态射的复合为 $(fg)_{\alpha}=\sum_{\alpha=\beta\gamma}f_{\beta}\beta_{\cdot}(g_{\gamma})$。这是加法范畴。置 $\mathrm{hot}_{-}(\mathcal{M}_{\cdot})=(\mathrm{hot}_{+}(\mathcal{M}^{\circ}_{\cdot}))^{\circ}$。有相应的复形 DG 范畴 $C\mathrm{hot}_{\pm}$。

有 DG 函子 $t_{+}:C\sec_{+}\to C\mathrm{hot}_{+}$，它把 $M\in C\sec_{+}$ 送到复形 $t_{+}M\in C\mathrm{hot}_{+}$，其分量为 $(t_{+}M)^{a}_{m}=M^{a-m}_{m}$，而微分 $d=d_{t_{+}M}$ 满足
\[
d_{\mathrm{id}_{\Delta_{m}}}=(-1)^{m}d^{a-m}_{M_{m}}:M^{a-m}_{m}\to M^{a-m+1}_{m},
\]
且对第 $i$ 个边缘映射 $\alpha_{i}:\Delta_{m}\hookrightarrow\Delta_{m+1}$ 有
\[
d_{\alpha_{i}}=(-1)^{i}\alpha^{*}_{i}:\alpha_{i\cdot}M^{a-m}_{m}\to M^{a-m}_{m+1},
\]
$d$ 的其余分量全为零。对 $l\in Hom(M_{1},M_{2})$ 有 $t_{+}(l)_{\mathrm{id}_{\Delta_{m}}}=l_{m}$，其余分量全为零。

\emph{注.} 函子 $t_{+}$ 是忠实的。可以把 $C\mathrm{hot}_{+}$ 的对象视为 $\sec_{+}$ 中带额外高阶同伦的「广义复形」。

另有 DG 函子 $s_{-}:C\mathrm{hot}_{+}\to C\sec_{-}$，定义如下。对 $A\in C\mathrm{hot}_{+}$，复形 $s_{-}A$ 的分量为
\[
(s_{-}A)^{a}_{m}=\sum_{\beta:\Delta_{n}\to\Delta_{m}}\beta_{\cdot}A^{a}_{n}.
\]
对 $\alpha:\Delta_{m}\to\Delta_{l}$，相容态射 $\alpha_{*}:(s_{-}A)^{a}_{l}\to\alpha_{\cdot}(s_{-}A)^{a}_{m}$ 的分量 $\gamma_{\cdot}A^{a}_{k}\to\alpha_{\cdot}\beta_{\cdot}A^{a}_{n}$ 当 $k=n$、$\gamma=\alpha\beta$ 时等于 $\mathrm{id}_{\gamma_{\cdot}A^{a}_{k}}$，否则为零。微分 $d_{s_{-}A}:(s_{-}A)^{a}_{m}\to(s_{-}A)^{a+1}_{m}$ 的分量 $\gamma_{\cdot}A^{a}_{k}\to\alpha_{\cdot}\beta_{\cdot}A^{a}_{n}$ 当 $\beta=\gamma\delta$ 时等于 $\gamma_{\cdot}(d_{A\delta})$，否则为零。

\emph{注.} DG 函子 $s_{-}$ 是全忠实的。
